\chapter*{Abstract}

\begin{spacing}{1.0}
Network intrusion detection relies on signature-based systems, which miss attacks without a known
signature, and on data-driven detectors, which are brittle when traffic changes. Neurosymbolic
approaches promise to join the two by bringing domain knowledge into learning. This thesis builds a
benchmark that compares Logic Tensor Networks with symbolic, classical machine-learning and
deep-learning baselines on NSL-KDD, CIC-IDS2017 and CIC-DDoS2019, with the same network as the
purely neural counterparts and the same tuning budget. It measures detection quality, cost and
usability, also under stress: scarce labels, rare classes, unseen attacks, cross-dataset transfer,
degraded knowledge, noise and evasion. Two rule-level metrics complete it.

On these datasets and knowledge bases most results are negative. Domain knowledge used as a
constraint in training gives no demonstrated gain: its effect on the macro-F1 is $+0.008$, $+0.015$
and $+0.002$, significant on no dataset after Holm correction. It does not help with few labels
either: where the effect stands out of the seeds it is negative, at most 1.5 points. Rules
extracted from the data do not hold outside training: on the two datasets with a separate test part
their precision falls to 0.130 and 0.547, while the mechanism rules of the same bases keep 0.950 to
0.984. Mechanism rules matter at inference in one case, the transfer from CIC-IDS2017 to
CIC-DDoS2019: the learned models reach a mean macro-F1 of at most 0.4812, while the UDP-flood rule of the
target dataset, read alone, reaches 0.918; in training the same rule only lifts the logic-trained
network to the level of the best neural baseline. Within each dataset, letting the full knowledge base override the
predictions lowers every learned model. At inference a network trained with logic costs the same as
a plain one, and each axiom adds about 2 to 4\,\% to the time of a training epoch. On CIC-IDS2017 both
tree ensembles beat every network. The thesis closes with a decision checklist for when a
neurosymbolic approach is worth using and when it is not.
% FIXME(open): E14 and E15 have no results; when they arrive, add the streaming pipeline (C4) and its latency and throughput here.
\end{spacing}
